%-*-latex-*-
%-*-latex-*-
\newcommand\COURSE{ciss245}
\newcommand\ASSESSMENT{q7101}
\newcommand\ASSESSMENTTYPE{Quiz}
\newcommand\POINTS{\textwhite{xxx/xxx}}

\input{myquizpreamble}
\input{yliow}
\input{\COURSE}
%-*-latex-*-

%-*-latex-*-

%-*-latex-*-
\renewcommand\TITLE{\ASSESSMENTTYPE \ \ASSESSMENT}

\renewcommand\EMAIL{}
\renewcommand\AUTHOR{eagraham1@cougars.ccis.edu}

\textwidth=6in
\begin{document}
\topmatterthree

%------------------------------------------------------------------------------
\nextq
The following program uses the \verb!vec2d! library.
\begin{Verbatim}[frame=single, fontsize=\footnotesize]
#include <iostream>
#include "vec2d.h"

int main()
{
    vec2d v(1.2, 3.4);

    int i = 2;
    try
    {
        std::cout << v[0] << std::endl; // prints 1.2
        std::cout << v[1] << std::endl; // prints 3.4
        std::cout << v[i] << std::endl; // exception thrown if i < 0 or i > 1
    }
    catch (IndexError & e)
    {
        std::cout << "IndexError caught" << std::endl; // prints IndexError caught
    }
    return 0;
}
\end{Verbatim}
Modify the \verb!vec2d.h! file
so that the exception catching code runs when the value of \verb!i!
is not \verb!0! or \verb!1!.
\\
\ANSWER
\begin{answercode}
// file: vec2d.h
#ifndef VEC2D_H
#define VEC2D_H

class IndexError
{};

class vec2d
{
public:
    vec2d(double x, double y)
        : x_(x), y_(y)
    {}
    double operator[](int i) const
    {
        if (i > 1 || i < 0) throw IndexError();
        return (i == 0 ? x_ : y_);
    }
    // other methods not shown
private:
    double x_, y_;
};
\end{answercode}

%------------------------------------------------------------------------------
\nextq
Modify the given \verb!Rational.h! so that exceptions are thrown whenever
there is a denominator of zero.
\\
\ANSWER
\begin{answercode}
// file: Rational.h
#ifndef RATIONAL_H
#define RATIONAL_H

class ZeroDenominatorError
{};

class Rational
{
public:
    Rational(int n, int d)
        : n_(x), d_(d)
    {}
    const Rational & operator/=(const Rational & f) const
    {
        if (f.d == 0 || d_ == 0) throw ZeroDenominatorError(); 
        n_ /= f.d_;
        d_ /= f.n_;
        return (*this);
    }
    Rational operator/(const Rational & f) const
    {
        if (f.d == 0 || d_ == 0) throw ZeroDenominatorError(); 
        return (Rational(*this) /= f);
    }
private:
    int n_, d_;
};
\end{answercode}
(Hint: Modify at two places.)

%------------------------------------------------------------------------------
%-*-latex-*-
\newpage
\input{instructions}
\end{document}

